\documentclass[10pt]{article}
\usepackage[T1]{fontenc}
\usepackage[utf8]{inputenc}
\usepackage[french]{babel}
\usepackage{amsmath,amssymb,mathrsfs,mathpazo}
\usepackage{geometry}
\geometry{a4paper,margin=2cm}
\usepackage{graphicx,array,tabularx,tabularray,colortbl,multirow,multicol,booktabs,enumitem,pifont,ulem,fancybox,diagbox,cancel}
\usepackage{pstricks,pstricks-add,pst-plot,pst-tree,pst-node,pst-text,pst-func,pst-math,pst-bspline,pst-3dplot,pst-eucl,pst-solides3d}
\let\pscalculate\undefined
\usepackage{pst-bezier}
\usepackage{tikz,pgfplots,tkz-tab}
\usetikzlibrary{arrows,arrows.meta,patterns,calc,positioning,trees,decorations.markings,shapes.geometric,shapes.misc,shapes.arrows,intersections,decorations.pathmorphing,decorations.pathreplacing,backgrounds,shadows,babel}
\pgfplotsset{compat=1.18}
\usepgfplotslibrary{fillbetween}
\usepackage[np]{numprint}
\usepackage{listings,makecell,pas-tableur,twoopt}
\let\barre\undefined
\usepackage[pstricks]{tablvar}
\definecolor{darkgreen}{RGB}{0,100,0}
\providecommand{\red}{\color{red}}
\providecommand{\blue}{\color{blue}}
\providecommand{\green}{\color{green}}
\providecommand{\black}{\color{black}}
\providecommand{\white}{\color{white}}
\providecommand{\gray}{\color{gray}}
\providecommand{\cyan}{\color{cyan}}
\providecommand{\magenta}{\color{magenta}}
\providecommand{\yellow}{\color{yellow}}
\providecommand{\violet}{\color{violet}}
\providecommand{\orange}{\color{orange}}
\providecommand{\euro}{\texteuro}
\providecommand{\EUR}{\texteuro}
\providecommand{\ieme}{\textsuperscript{e}}
\providecommand{\degres}{\ensuremath{^\circ}}
\providecommand{\degree}{\ensuremath{^\circ}}
\providecommand{\vv}[1]{\overrightarrow{#1}}
\providecommand{\Coord}[1]{#1}
\providecommand{\decosix}{}
\providecommand{\starredbullet}{\ensuremath{\bullet}}
\pagestyle{empty}
\setlength{\parindent}{0pt}
\usepackage{adjustbox}\DeclareUnicodeCharacter{2236}{:}\DeclareUnicodeCharacter{2190}{\ensuremath{\leftarrow}}
\providecommand{\R}{\mathbb{R}}
\providecommand{\N}{\mathbb{N}}
\providecommand{\Z}{\mathbb{Z}}
\providecommand{\Q}{\mathbb{Q}}
\providecommand{\D}{\mathbb{D}}
\providecommand{\C}{\mathbb{C}}
\providecommand{\e}{{\mathrm{e}}}
\providecommand{\vect}[1]{\overrightarrow{#1}}
\providecommand{\vectt}[1]{\overrightarrow{\mathrm{#1}}}
\providecommand{\vvt}[1]{\overrightarrow{\mathrm{#1}}}
\providecommand{\np}[2][]{{#2~#1}}
\providecommand{\vv}[1]{{\overrightarrow{#1}}}
\renewcommand{\d}{{\mathrm{d}}}
\providecommand{\degres}{{^\circ}}
\providecommand{\degre}{{^\circ}}
\providecommand{\blue}{\color{blue}}
\providecommand{\red}{\color{red}}
\providecommand{\enskip}{\quad}
\providecommand{\llbracket}{\lbrack\!\lbrack}
\providecommand{\rrbracket}{\rbrack\!\rbrack}
\providecommand{\euro}{{\text{€}}}
\providecommand{\no}{{\text{n°}}}
\providecommand{\interval}[2]{{[#1;#2]}}
\providecommand{\intervFO}[2]{{[#1;#2[}}
\providecommand{\dotfill}{\ldots}

\providecommand{\ds}{\displaystyle}
\providecommand{\textding}[1]{#1}
\providecommand{\Oij}{\(\mathrm{O};\vec{\imath},\vec{\jmath}\)}
\providecommand{\Oijk}{\(\mathrm{O};\vec{\imath},\vec{\jmath},\vec{k}\)}
\providecommand{\Ouv}{\(\mathrm{O};\vec{u},\vec{v}\)}
\providecommand{\intervFF}[2]{\ensuremath{[#1;#2]}}
\renewcommand{\euro}{\ifmmode\text{\texteuro}\else\texteuro\fi}
\DeclareUnicodeCharacter{00B0}{\ifmmode^{\circ}\else\textdegree\fi}
\DeclareUnicodeCharacter{20AC}{\ifmmode\text{\texteuro}\else\texteuro\fi}
% Énoncés : Éducation nationale. Transcription et corrigés : APMEP, Denis Vergès et contributeurs. Découpage et classement : CoopMaths, Sébastien Lozano, Éric Elter et contributeurs. Adaptation de la présentation : Maths Explorer.
% CC BY-SA : https://coopmaths.fr/www/about/
% https://coopmaths.fr/alea/static/bac/2023/tex/bac_2023_09_sujet2_ameriquesud_1.tex
\begin{document}
\noindent\textit{APMEP / CoopMaths. CC BY-SA. Présentation adaptée par Maths Explorer.}\par\medskip
\medskip

\textbf{Partie A}

\medskip

Un jeu proposé dans une fête foraine consiste à effectuer trois tirs successivement sur une cible mouvante.

On a constaté que :
\begin{itemize}
\item[$\bullet~~$] Si le joueur atteint la cible lors d'un tir alors il ne l'atteint pas lors du tir suivant dans 65\,\% des cas ;
\item[$\bullet~~$] Si le joueur n'atteint pas la cible lors d'un tir alors il l'atteint lors du tir suivant dans 50\,\% des cas.
\end{itemize}

La probabilité qu'un joueur atteigne la cible lors de son premier tir est de $0,6$.

\emph{Pour tout évènement $A$, on note $p(A)$ sa probabilité et $\overline{A}$ l'évènement contraire de $A$.}

On choisit au hasard un joueur à ce jeu de tirs.

On considère les évènements suivants:
\begin{itemize}
\item[$\bullet~~$] $A_1$ : \og Le joueur atteint la cible lors du 1\up{er} tir \fg
\item[$\bullet~~$] $A_2$ : \og Le joueur atteint la cible lors du 2\up{e} tir \fg
\item[$\bullet~~$] $A_3$ : \og Le joueur atteint la cible lors du 3\up{e} tir \fg.
\end{itemize}

\medskip

\begin{enumerate}
\item Recopier et compléter, avec les probabilités correspondantes sur chaque branche, l'arbre pondéré ci-dessous modélisant la situation.

\begin{center}
\pstree[treemode=R,nodesepA=0pt,nodesepB=2.5pt,treesep = 0.75cm,levelsep=3cm]{\TR{}}
{\pstree{\TR{$A_1$~~}\taput{$0,6$}}
	{\pstree{\TR{$A_{2}$~~}}
		{
		\TR{$A_{3}~~$}
		\TR{$\overline{A_{3}~~}$}
		}
	\pstree{\TR{$\overline{A_{2}~~}$}}
		{
		\TR{$A_{3}~~$}
		\TR{$\overline{A_{3}~~}$}
		}
	}
\pstree{\TR{$\overline{A_1}$~~}}
	{\pstree{\TR{$A_{2}$~~}}
		{
		\TR{$A_{3}~~$}
		\TR{$\overline{A_{3}~~}$}
		}
	\pstree{\TR{$\overline{A_{2}}$~~}}
		{
		\TR{$A_{3}$~~}
		\TR{$\overline{A_{3}~~}$}
		}
	}
}
\end{center}

Soit $X$ la variable aléatoire qui donne le nombre de fois où le joueur atteint sa cible au cours
des trois tirs.
\item Montrer que la probabilité que le joueur atteigne exactement deux fois la cible au cours des trois tirs est égale à \np{0,4015}.
\item L'objectif de cette question est de calculer l'espérance de la variable aléatoire $X$, notée $E(X)$.
	\begin{enumerate}
		\item Recopier et compléter le tableau ci-dessous donnant la loi de probabilité de la variable aléatoire $X$.
\begin{center}
\begin{tabularx}{\linewidth}{|*{5}{>{\centering \arraybackslash}X|}}\hline
$X = x_i$				&0	&1	&2	&3\\ \hline
$p\left(X = x_i\right)$	&0,1&	&	& \np{0,0735}\\ \hline
\end{tabularx}
\end{center}
		\item Calculer $E(X)$.
		\item Interpréter le résultat précédent dans le contexte de l'exercice.
	\end{enumerate}
\end{enumerate}

\medskip

\textbf{Partie B}

\medskip

On considère $N$, un entier naturel supérieur ou égal à 1.

Un groupe de $N$ personnes se présente à ce stand pour jouer à ce jeu dans des conditions identiques et indépendantes. 

Un joueur est déclaré gagnant lorsqu'il atteint trois fois la cible.

On note $Y$ la variable aléatoire qui compte parmi les $N$ personnes le nombre de joueurs déclarés gagnants.

\medskip

\begin{enumerate}
\item Dans cette question, $N = 15$.
	\begin{enumerate}
		\item Justifier que $Y$ suit une loi binomiale dont on déterminera les paramètres.
		\item Donner la probabilité, arrondie à $10^{-3}$, qu'exactement 5 joueurs soient gagnants à ce jeu.
	\end{enumerate}
\item Par la méthode de votre choix, que vous expliciterez, déterminer le nombre minimum de personnes qui doivent se présenter à ce stand pour que la probabilité qu'il y ait au moins un joueur gagnant soit supérieure ou égale à $0,98$.
\end{enumerate}

\bigskip

\end{document}
